\section{Report}
\label{sec:report}

We report results under an audit-first protocol that separates (i) outcome metrics from (ii) process and
format validity.

\paragraph{Three-tier denominators (T/I/M).}
To avoid silent pruning when a metric is undefined, we disclose data under three nested tiers:
(T) \emph{Total} includes all collected submissions and reports rates of OOS, missingness (NA), and other
process events; (I) \emph{In-scope} restricts to tasks with non-missing \texttt{scope} marked in-scope;
(M) \emph{Model-usable} further excludes pre-registered, metric-specific gating failures.
All exclusions are disclosed with explicit counts and reasons.

\paragraph{Type~4 (process/format) disclosures.}
We distinguish system-side NA (e.g., missing/corrupted exports) from non-compliant multi-select meta-label
records (empty selections when required; mutually-exclusive co-selection such as \texttt{trivial}+others or
\texttt{acceptable}+others). UI-level hard validation aims to drive empty/conflict rates near zero, but NA
cannot be fully eliminated; therefore both are disclosed with explicit denominators.

\paragraph{Process evidence for ``blocked errors''.}
When submission-time validation prevents non-compliant records from entering the main dataset, we disclose
the number of blocked submission attempts and rejection reasons (e.g., empty vs. mutual-exclusion conflict),
optionally stratified by worker, as a process-quality indicator.

\paragraph{Worker profiling outputs.}
Worker-level reporting uses $\mathrm{LCB}(r_w)$ together with the discrete risk tier $R_w^{\mathrm{tier}}$ rather
than a continuous anomaly score. Report tables and figures additionally disclose the dominant risk signature
(blind trust, condition fragility, or gate-failure tendency) and the rule version/reason code used to assign
each worker to a tier.
